\documentclass[coursework]{SCWorks}
\usepackage{preamble}
\setminted[py]{fontsize=\small, breaklines=true, style=bw, linenos}
\newcommand{\eqdef}{\stackrel {\rm def}{=}}
\newcommand{\specialcell}[2][c]{%
    \begin{tabular}[#1]{@{}c@{}}#2\end{tabular}}

\renewcommand\theFancyVerbLine{\small\arabic{FancyVerbLine}}

\newtheorem{lem}{Лемма}
\linespread{1.5}
\setminted{
    fontsize=\small,          % Размер шрифта
    baselinestretch=1.0,      % Единичный интервал
    breaklines=true,          % Перенос строк
    framesep=2mm,             % Отступ вокруг кода
    linenos=false,            % Нумерация строк (опционально)
    xleftmargin=2em,          % Отступ слева
    breakanywhere=true,
}

\begin{document}

    Работа посвящена сравнительному анализу моделей генерации заголовков для русского языка.

    \subsection{Сбор и предобработка данных}

    Собран корпус из 30 000 пар «текст — заголовок».
    Для оптимизации обучения тексты разбивались на фрагменты по 2–3 предложения.
    Для исключения переобучения введено ограничение на количество повторяющихся заголовков.

    \subsection{Модели}

    \subsubsection{N-gram модель}

    N-gram модель строит вероятностное распределение слов заголовка на основе контекста из n предыдущих слов.
    Использовалась модель 4-го порядка со сглаживанием Лапласа и стратегией backoff

    N-gram модель продемонстрировала возможность решения задачи, однако страдает от ограничений
    статистического подхода: неспособность к обобщению и зависимость от точного совпадения контекста.

    \subsubsection{Полносвязная языковая модель (FNNLM)}

    Feed-Forward Neural Language Model — модели с фиксированным окном контекста в 8 токенов.

    Несмотря на лучшую обобщающую способность, фиксированное окно ограничивает учет глобального контекста.

    \subsubsection{Рекуррентная нейронная сеть (RNN)}

    Архитектура работает в два этапа: сначала текст кодируется через RNN для получения финального скрытого состояния, затем
    на основе этого состояния последовательно генерируются токены заголовка.

    Ключевое ограничение данной архитектуры — проблема затухания градиента: RNN неспособна эффективно
    запоминать долгосрочные зависимости.

    \subsubsection{Двунаправленная модель долгой краткосрочной памяти (BiLSTM)}

    Модель LSTM с двунаправленным кодированием решила проблему долгосрочных зависимостей.

    Текст и заголовок объединяются в единую последовательность с разделительным токеном, а функция потерь
    вычисляется только по части заголовка благодаря маске.

    Двунаправленное кодирование позволяет модели учитывать контекст, что существенно улучшает качество генерации по сравнению
    с однонаправленной RNN.

    \subsubsection{Seq2Seq с механизмом внимания}

    Архитектура Seq2Seq с механизмом внимания представляет собой наиболее продвинутый подход среди реализованных и состоит
    из трёх основных компонентов.

    \begin{itemize}

        \item \textbf{Encoder:} Двунаправленный LSTM, формирующий сжатое представление текста.

        \item \textbf{Attention:} Позволяет декодеру фокусироваться на релевантных частях входа при генерации каждого токена.

        \item \textbf{Decoder:} LSTM, использующий вектор внимания для формирования итоговой последовательности.

    \end{itemize}

    Модель превосходит базовые за счет способности «выхватывать» важные образы из любой части текста.

    \subsection{Заключение}

    Результаты подтверждают, что для русского языка критически важны механизмы внимания и двунаправленного кодирования.
    Перспективным направлением является переход к трансформерным архитектурам (BERT/GPT-like)
    для работы с ограниченными объемами специфических данных.
\end{document}